\documentclass[12pt]{article}
%^ Start of document
\usepackage{times}  %Makes text TimesNewRoman
\usepackage{setspace} %Allows you to set single or double space 
\usepackage{biblatex} %Imports biblatex package
\usepackage{graphicx} % Required for inserting images
\usepackage{authblk}
\usepackage[colorlinks=true, urlcolor=blue, linkcolor=blue]{hyperref}
\usepackage{subfig}
\usepackage{float}
\usepackage{tabularx}
\usepackage{multirow}
\usepackage[paperheight=11in,
   paperwidth=8.5in,
   top=20mm,
   bottom=20mm,
   left=40mm,
   right=40mm]{geometry}
\usepackage{moresize}
\usepackage{tikz}
\usepackage{xcolor}
\usepackage{physics}
\usepackage{amssymb}
\usepackage{mathrsfs}
\usepackage{amsmath}
\usepackage[document]{ragged2e}
\usepackage{comment}
\usepackage{xcolor}
\addbibresource{Qhw.bib}
\begin{document}
\singlespacing  %Sets single spacing for whole document
\title{Astrophysics Midterm}

\author[1]{Zachary Tullier}
\affil[1]{Student\\Physics Department\\ University of Houston - Clear Lake \\Houston, Texas\\U.S}
\date{}
\maketitle

\section*{Problem 1: Hot Star in M31}

A star in the Andromeda galaxy (M31) is observed to have a flux
\[
F = 10^{-12}~\mathrm{erg~cm^{-2}~s^{-1}},
\]
and is assumed to radiate as a blackbody with temperature
\[
T = 6\times 10^{5}~\mathrm{K}.
\]
The distance to M31 is taken as
\[
d = 800~\mathrm{kpc}.
\]

\subsection*{(a) Luminosity}

Convert distance to centimeters:
\[
d = 800~\mathrm{kpc}
  = 800\times 10^{3}~\mathrm{pc}
  = 8\times 10^{5}~\mathrm{pc}
  = 8\times 10^{5}(3.086\times 10^{18}~\mathrm{cm})
  = 2.47\times 10^{24}~\mathrm{cm}.
\]
The luminosity is
\[
L = 4\pi d^{2}F.
\]
Compute:
\[
d^{2} = (2.47\times 10^{24})^{2} = 6.10\times 10^{48}~\mathrm{cm^{2}},
\]
\[
L = 4\pi(6.10\times 10^{48})(10^{-12})
  = 7.7\times 10^{37}~\mathrm{erg~s^{-1}}.
\]
\[
\boxed{L \approx 7.7\times 10^{37}~\mathrm{erg~s^{-1}}}
\]

\subsection*{(b) Radius}

Using the Stefan--Boltzmann law:
\[
L = 4\pi R^{2}\sigma T^{4},
\qquad
R = \sqrt{\frac{L}{4\pi\sigma T^{4}}}.
\]
Compute:
\[
T^{4} = (6\times 10^{5})^{4}
      = 1.296\times 10^{23},
\]
\[
4\pi\sigma T^{4}
  = 12.57(5.67\times 10^{-5})(1.296\times 10^{23})
  = 9.24\times 10^{19}.
\]
Thus,
\[
R^{2} = \frac{7.7\times 10^{37}}{9.24\times 10^{19}}
      = 8.3\times 10^{17},
\]
\[
R = 9.1\times 10^{8}~\mathrm{cm} = 9.1\times 10^{3}~\mathrm{km}.
\]
\[
\boxed{R \approx 9\times 10^{8}~\mathrm{cm}}
\]

\subsection*{(c) Peak Wavelength}

Using Wien's law:
\[
\lambda_{\max} = \frac{b}{T}
= \frac{2.898\times 10^{-3}}{6\times 10^{5}}
= 4.83\times 10^{-9}~\mathrm{m}
= 4.8~\mathrm{nm}
= 48~\text{\AA}.
\]
\[
\boxed{\lambda_{\max} \approx 4.8~\mathrm{nm}}
\]

\section*{Problem 2: Binary System with an X-ray Pulsar}

The system has an orbital period
\[
P = 10~\mathrm{days} = 8.64\times 10^{5}~\mathrm{s},
\]
optical star velocity
\[
v_{1} = 20~\mathrm{km\,s^{-1}} = 2\times 10^{4}~\mathrm{m\,s^{-1}},
\]
and pulsar orbital radius
\[
r_{2} = 3\times 10^{12}~\mathrm{cm} = 3\times 10^{10}~\mathrm{m}.
\]

\subsection*{(a) Orbit radius of star 1}

Circular orbit:
\[
v_{1} = \frac{2\pi r_{1}}{P}
\Rightarrow
r_{1} = \frac{v_{1}P}{2\pi}.
\]
\[
r_{1} = \frac{(2\times 10^{4})(8.64\times 10^{5})}{6.283}
      = 2.75\times 10^{9}~\mathrm{m}.
\]
\[
\boxed{r_{1} \approx 2.8\times 10^{9}~\mathrm{m}}
\]

\subsection*{(b) Separation}

\[
a = r_{1} + r_{2}
= 2.75\times 10^{9} + 3\times 10^{10}
= 3.28\times 10^{10}~\mathrm{m}.
\]
\[
\boxed{a \approx 3.3\times 10^{10}~\mathrm{m}}
\]

\subsection*{(c) Total Mass}

Newtonian form of Kepler’s third law:
\[
M_{\mathrm{tot}} = \frac{4\pi^{2}a^{3}}{GP^{2}}.
\]
Compute:
\[
a^{3} = (3.28\times 10^{10})^{3}
      = 3.53\times 10^{31}~\mathrm{m^{3}},
\]
\[
4\pi^{2}a^{3}
= 1.39\times 10^{33},
\]
\[
GP^{2}
= (6.674\times 10^{-11})(7.47\times 10^{11})
= 49.9.
\]
Thus:
\[
M_{\mathrm{tot}}
= \frac{1.39\times 10^{33}}{49.9}
= 2.8\times 10^{31}~\mathrm{kg}
\approx 14\,M_{\odot}.
\]
\[
\boxed{M_{\mathrm{tot}} \approx 14\,M_{\odot}}
\]

\subsection*{(d) Individual Masses}

Center-of-mass condition:
\[
M_{1}r_{1} = M_{2}r_{2},
\qquad
\frac{M_{2}}{M_{1}} = \frac{r_{1}}{r_{2}}
= 0.0917.
\]
Let $q = M_{2}/M_{1}$. Then
\[
M_{1} = \frac{M_{\mathrm{tot}}}{1+q}
= \frac{2.8\times 10^{31}}{1.0917}
= 2.55\times 10^{31}~\mathrm{kg}
\approx 12.8\,M_{\odot},
\]
\[
M_{2} = M_{\mathrm{tot}} - M_{1}
= 2.3\times 10^{30}~\mathrm{kg}
\approx 1.2\,M_{\odot}.
\]
\[
\boxed{M_{1} \approx 13\,M_{\odot}},\qquad
\boxed{M_{2} \approx 1.2\,M_{\odot}}
\]

\section*{Problem 3: The Standard Solar Model}

The \emph{Standard Solar Model} (SSM) is a self-consistent description of the Sun’s
interior and evolution based on hydrostatic equilibrium, nuclear fusion,
radiative/convective energy transport, and realistic microphysics.

\subsection*{Core}

The core ($0$--$0.25R_{\odot}$) contains the temperature and density required for
hydrogen fusion:
\[
T_{\mathrm{c}}\sim 1.5\times 10^{7}~\mathrm{K},\qquad
\rho_{\mathrm{c}}\sim 150~\mathrm{g\,cm^{-3}}.
\]
Energy is generated through the proton--proton chain and, at a lower level, the
CNO cycle:
\[
p + p \rightarrow d + e^{+} + \nu_{e},\qquad
^{3}\mathrm{He} + ^{3}\mathrm{He} \rightarrow \, ^{4}\mathrm{He} + 2p.
\]
Most of the Sun’s luminosity arises from this region.

\subsection*{Radiative Zone}

From $0.25R_{\odot}$ to $0.7R_{\odot}$ energy is transported by radiation.  The
temperature decreases smoothly from several million kelvin.  Photons undergo a
random walk, giving diffusion times of $10^{5}$--$10^{6}$ years.

\subsection*{Convective Zone}

Above $0.7R_{\odot}$, rising opacity and partial ionization of H and He make
radiative transport inefficient.  The layer becomes convectively unstable.
Energy transport is dominated by bulk motions.  This convection drives
photospheric granulation.

\subsection*{Photosphere}

The visible ``surface’’ of the Sun, with
\[
T_{\mathrm{eff}} \approx 5770~\mathrm{K}.
\]
It produces the continuous optical spectrum plus Fraunhofer absorption lines.
Granulation is observed here.

\subsection*{Chromosphere}

A region of low density where the temperature rises from $4500~\mathrm{K}$ to
$\sim 10^{4}~\mathrm{K}$.  Visible in H$\alpha$ as a reddish rim during eclipses.
Contains spicules and is governed strongly by magnetic fields.

\subsection*{Transition Region}

A thin layer in which temperature rises sharply by orders of magnitude, linking
the chromosphere to the corona.

\subsection*{Corona}

A hot, tenuous, magnetically structured plasma with
\[
T \sim 10^{6}~\mathrm{K}.
\]
It emits X-rays and EUV radiation.  The mechanism of coronal heating remains an
active research topic.

\subsection*{Solar Wind}

The corona extends outward into interplanetary space as a supersonic plasma
flow, carrying magnetic fields into the heliosphere.

\section*{Conclusion}

The Standard Solar Model provides a consistent physical description of the Sun
from core to corona, matching observed luminosity, radius, helioseismic data,
and solar neutrino measurements (after accounting for neutrino oscillations).
\end{document}
